\begin{corollary}[Dirichlet L-function at positive integers for even/odd characters] \label{corollary:dirichlet_L_function_at_positive_integers_even_odd_character}
\TODO{AI generated, verify}
Let $\chi$ be a \CrefAndHyperrefIfExist{definition:dirichlet_character_modulo_a_nonnegative_integer}{Dirichlet character modulo $k$}. For $n \in \bbN$ with $n \equiv a(\chi) \pmod{2}$ (where $a(\chi) = 0$ if $\chi(-1)=1$ and $a(\chi)=1$ if $\chi(-1)=-1$), the value $L(n, \chi)$ is a non-zero algebraic multiple of $\pi^n$ times an appropriate Gauss sum. In particular, for even characters ($a=0$), $L(2n, \chi) \in \overline{\bbQ} \cdot \pi^{2n}$; for odd characters ($a=1$), $L(2n+1, \chi) \in \overline{\bbQ} \cdot \pi^{2n+1}$.

This follows from the functional equation in \Cref{theorem:analytic_continuation_and_functional_equation_for_dirichlet_L_function} and the analytic class number formula for Dirichlet L-functions.

\CrefAndHyperrefIfExist{definition:dirichlet_character_modulo_a_nonnegative_integer}{Dirichlet character}
\CrefAndHyperrefIfExist{definition:dirichlet_L_function_of_a_dirichlet_character_modulo_a_nonnegative_integer}{Dirichlet L-function}
\Cref{theorem:analytic_continuation_and_functional_equation_for_dirichlet_L_function}
\end{corollary}